\documentclass[12pt]{article}

% Overleaf: set the compiler to pdfLaTeX and bibliography tool to Biber.
% This template uses APA author-year citations through biblatex.

\usepackage[letterpaper,margin=1in]{geometry}
\usepackage[T1]{fontenc}
\usepackage[utf8]{inputenc}
\usepackage{lmodern}
\usepackage{setspace}
\usepackage{booktabs}
\usepackage{array}
\usepackage{enumitem}
\usepackage{hyperref}
\usepackage{csquotes}
\usepackage[
  backend=biber,
  style=apa,
  sorting=nyt,
  sortcites=true
]{biblatex}

\addbibresource{references.bib}

\hypersetup{
  colorlinks=true,
  linkcolor=black,
  citecolor=black,
  urlcolor=blue
}

\setstretch{1.5}
\setlist[itemize]{leftmargin=*, itemsep=0.25em}
\setlist[enumerate]{leftmargin=*, itemsep=0.25em}

\title{Research Proposal Title}
\author{Your Name\\Department or Program\\Brock University}
\date{\today}

\begin{document}

\maketitle

\begin{abstract}
Write a concise summary of the proposed project. In 150--250 words, identify
the research problem, the central research question or hypothesis, the proposed
method, and the expected contribution. Keep this focused on what the project
will do, why it matters, and why it is feasible.
\end{abstract}

\section{Introduction}

The introduction orients readers who may not already know the topic. Begin with
a brief overview of the research area, then narrow toward the specific problem
your project addresses.

\subsection{Research Topic and Problem}

Describe the topic and define the specific research problem. Be selective:
include only the background needed to understand the proposal.

\subsection{Research Question or Hypothesis}

State the main research question, sub-question, or hypothesis clearly.

\begin{quote}
\textbf{Main research question:} Replace this sentence with one focused,
researchable question.
\end{quote}

\subsection{Objectives}

List the concrete objectives of the proposed project.

\begin{enumerate}
  \item Objective one.
  \item Objective two.
  \item Objective three.
\end{enumerate}

\subsection{Significance}

Explain the expected contribution, impact, or importance of the project. This
section should answer why this research should be done.

\section{Literature Review}

The literature review explains the scholarly conversation surrounding your
topic. It should support the need for your research question or hypothesis.
Avoid summarizing everything you have read; focus on the sources most relevant
to your project.

\subsection{Current Scholarly Conversation}

Summarize the main areas of agreement in the literature. Use citations such as
\textcite{sampleArticle} for narrative citation or \parencite{sampleBook} for
parenthetical citation.

\subsection{Debates, Conflicts, or Limitations}

Identify important disagreements, limitations, or unresolved questions in
existing research.

\subsection{Gap and Connection to This Project}

Explain the gap your project addresses and connect it directly to your research
question or hypothesis. This section should make clear why your proposed project
fits into the larger conversation.

\section{Methodology}

The methodology section explains how the project will answer the research
question or test the hypothesis. It should justify both the method of data
collection and the method of analysis.

\subsection{Research Design}

Describe the overall approach, such as qualitative, quantitative, mixed-methods,
experimental, archival, textual, survey-based, interview-based, or another
appropriate design.

\subsection{Data Collection}

Describe what data, texts, participants, observations, records, or materials
will be collected. Explain why this type of evidence is appropriate for the
research question.

\subsection{Data Analysis}

Describe how the collected material will be analyzed. Explain how this analysis
will help answer the research question or test the hypothesis.

\subsection{Justification of Method}

Explain why the selected method is the best feasible approach. Address:

\begin{itemize}
  \item \textbf{Viability:} The project can be completed within available time,
  resources, expertise, and budget.
  \item \textbf{Reliability:} The approach has a track record or can be carried
  out consistently.
  \item \textbf{Validity:} The method produces evidence appropriate to the
  research question.
\end{itemize}

\subsection{Ethical Considerations}

If applicable, describe consent, privacy, confidentiality, data storage, risk,
and research ethics board requirements. If ethics approval is not required,
briefly explain why.

\section{Timeline}

Work backward from the planned completion date. Adjust the milestones below to
match your program requirements and project scope.

\begin{table}[h]
\centering
\caption{Proposed research timeline}
\begin{tabular}{p{0.28\linewidth} p{0.25\linewidth} p{0.37\linewidth}}
\toprule
\textbf{Milestone} & \textbf{Target Date} & \textbf{Planned Work} \\
\midrule
Finalize research question & Month Year & Confirm scope with supervisor and revise proposal focus. \\
Complete literature review & Month Year & Search, read, annotate, and synthesize key sources. \\
Finalize methodology & Month Year & Confirm data collection and analysis plan. \\
Collect data or materials & Month Year & Gather sources, observations, interviews, survey responses, or other evidence. \\
Analyze data & Month Year & Apply selected analytical method and record findings. \\
Draft chapters or report & Month Year & Write and revise major sections. \\
Submit final project & Month Year & Complete final revisions, formatting, and submission. \\
\bottomrule
\end{tabular}
\end{table}

\section{Expected Contribution}

Briefly state what this project is expected to contribute to the field,
discipline, community, practice, or problem area. Connect the contribution back
to the research gap and significance described earlier.

\section{Conclusion}

End by restating the focused purpose of the project, the value of the research,
and the feasibility of the plan. Keep the conclusion brief.

\printbibliography

\end{document}
